% ============================================================
\section{Solver Infrastructure: Assembly, Projection, Backend, and Utilities}
\label{sec:solver-infra}
% ============================================================

This chapter documents the horizontal infrastructure layers that sit between
the high-level \texttt{HybridPowerSystem} model and the numerical solver
kernels: the \textbf{assembly} layer (flat solver-ready data), the
\textbf{projection} layer (canonical model expansion and bus merging), the
\textbf{backend} abstraction interfaces (linear, NLP, and QP solvers), and
the \textbf{thread pool} utility.

% ------------------------------------------------------------
\subsection{Assembly Layer}
\label{sec:assembly}
% ------------------------------------------------------------

The assembly layer (\texttt{include/hacdcpf/assembly/}) converts a
\texttt{HybridPowerSystem} into flat, pre-computed data structures ready for
direct consumption by numerical solvers.

\subsubsection{SolverData}
\label{sec:solver-data}

\texttt{assembly/solver\_data.hpp} --- namespace \texttt{hacdcpf::powerflow}.

\texttt{SolverData} is the canonical flat representation of a projected
network.  It owns all component vectors and the precomputed sparse matrices
used by every PF/OPF kernel.

\paragraph{Component vectors}
\begin{longtable}{ll}
\hline\hline
\textbf{Field} & \textbf{Description} \\
\hline\endhead\hline\endfoot
\texttt{ac\_buses} & AC bus vector (post-projection) \\
\texttt{ac\_branches} & AC branch vector (incl.\ expanded transformer equivalents) \\
\texttt{dc\_buses} & DC bus vector \\
\texttt{dc\_branches} & DC branch vector \\
\texttt{converters} & VSC AC/DC converters \\
\texttt{dcdc\_converters} & DC-DC converters \\
\texttt{energy\_routers} & Multi-port energy routers \\
\texttt{generators} & Synchronous and asynchronous generators \\
\texttt{loads} & AC constant-power and ZIP loads \\
\texttt{flexible\_loads} & Demand-response flexible loads \\
\texttt{static\_generators} & Distributed generation (inverter-coupled) \\
\texttt{renewable\_gens} & Wind and solar renewable generation \\
\texttt{pv\_systems} & PV array systems (MPPT model) \\
\texttt{storage\_units} & AC battery / flywheel storage \\
\texttt{dc\_loads} & DC bus loads \\
\texttt{dc\_storage} & DC bus storage \\
\texttt{dc\_static\_generators} & DC-connected distributed generation \\
\texttt{dc\_pv\_arrays} & DC PV arrays \\
\texttt{shunts} & Fixed and switchable shunts \\
\texttt{charging\_stations} & EV charging stations \\
\texttt{chargers} & Individual EV chargers \\
\texttt{external\_grids} & External grid equivalents \\
\texttt{switches} & Switching devices (CB, DS, LBS, etc.) \\
\texttt{vpps} & Virtual power plants \\
\texttt{microgrids} & Microgrid aggregates \\
\texttt{mobile\_storage} & Mobile storage units \\
\end{longtable}

\paragraph{Numerical parameters and flags}
\begin{longtable}{lll}
\hline\hline
\textbf{Field} & \textbf{Default} & \textbf{Description} \\
\hline\endhead\hline\endfoot
\texttt{base\_mva} & 100.0 & System base MVA \\
\texttt{loss\_model} & \texttt{Linear} & VSC converter loss model type \\
\texttt{zip\_pw[3]} & \{1,0,0\} & ZIP P weights (Z,I,P fractions) \\
\texttt{zip\_qw[3]} & \{1,0,0\} & ZIP Q weights \\
\texttt{has\_component\_loads} & false & True after ZIP load aggregation \\
\texttt{pd\_pu}, \texttt{qd\_pu} & --- & Aggregated bus P/Q load (pu) \\
\texttt{bus\_zip\_*} & --- & Per-bus ZIP weight vectors (Z/I/P, P and Q) \\
\texttt{enable\_coupled\_jacobian} & false & Coupled AC/DC Jacobian \\
\texttt{enable\_augmented\_equations} & false & VDC/VAC augmented rows \\
\texttt{enable\_semi\_smooth\_newton} & false & Fischer-Burmeister PV/PQ \\
\texttt{ncp\_mu} & 0.0 & Smooth-NCP smoothing parameter \\
\texttt{pg}, \texttt{qg} & --- & Aggregated bus generator P/Q (pu) \\
\texttt{ybus} & --- & Complex admittance matrix \(\mathbf{Y}_{bus}\) (sparse) \\
\texttt{gdc} & --- & DC conductance matrix \(\mathbf{G}_{dc}\) (sparse) \\
\texttt{build\_id} & 0 & Monotone counter incremented on each rebuild \\
\texttt{bus\_merge\_map} & nullopt & Optional zero-impedance bus merge map \\
\end{longtable}

\paragraph{Factory and utility functions}
\begin{center}
\begin{tabular}{ll}
\hline\hline
\textbf{Function} & \textbf{Description} \\
\hline
\texttt{make\_solver\_data(sys)} & Build from \texttt{HybridPowerSystem} (copy) \\
\texttt{make\_solver\_data(sys\&\&)} & Build from rvalue (move-efficient) \\
\texttt{make\_solver\_data\_projected(projected\&\&)} & Build from already-projected system \\
\texttt{aggregate\_generation(data)} & Sum generator injections onto bus vectors \\
\texttt{aggregate\_load\_demand(data)} & Sum load demands onto bus vectors \\
\texttt{rebuild\_matrices(data)} & Recompute \(\mathbf{Y}_{bus}\) and \(\mathbf{G}_{dc}\) \\
\hline\hline
\end{tabular}
\end{center}

\subsubsection{Index Maps}
\label{sec:index-map}

\texttt{assembly/index\_map.hpp} maps between external (1-based user-facing)
and internal (0-based solver) bus and branch indices.

\begin{center}
\begin{tabular}{ll}
\hline\hline
\textbf{Type} & \textbf{Description} \\
\hline
\texttt{BusIndexMap::ext\_to\_pos} & \texttt{unordered\_map< int,int>}: external $\to$ internal \\
\texttt{BusIndexMap::pos\_to\_ext} & \texttt{vector< int>}: internal $\to$ external \\
\texttt{BranchIndexMap} & Same pattern for branches \\
\texttt{SystemIndexMap} & Combined AC/DC bus and branch maps + optional merge map pointer \\
\hline\hline
\end{tabular}
\end{center}

\subsubsection{Admittance (Y-bus) Builder}
\label{sec:ybus-builder}

\texttt{assembly/ybus\_builder.hpp} provides three matrix builders:

\paragraph{Functions}
\begin{center}
\begin{tabular}{ll}
\hline\hline
\textbf{Function} & \textbf{Returns} \\
\hline
\texttt{build\_admittance\_matrix(data)} & Full complex \(\mathbf{Y}_{bus}\) for AC Newton PF \\
\texttt{build\_dc\_conductance(data)} & Real \(\mathbf{G}_{dc}\) for DC PF \\
\texttt{build\_susceptance\_matrix(data, flags)} & Real \(\mathbf{B}'\) or \(\mathbf{B}''\) for FDPF \\
\hline\hline
\end{tabular}
\end{center}

\paragraph{BpBuildFlags} (used with \texttt{build\_susceptance\_matrix}):
\begin{longtable}{lll}
\hline\hline
\textbf{Flag} & \textbf{Default} & \textbf{Effect} \\
\hline\endhead\hline\endfoot
\texttt{zero\_resistance} & false & Set series $R=0$ (builds $B'$ for XB-FDPF) \\
\texttt{zero\_charging} & false & Ignore line charging $b_{pu}$ \\
\texttt{unit\_tap} & false & Treat tap magnitude as 1 \\
\texttt{zero\_phase\_shift} & false & Zero all transformer phase shifts \\
\texttt{add\_bus\_shunts} & false & Add bus shunt susceptance $-B_s$ (builds $B''$) \\
\texttt{min\_x\_pu} & 0.0 & Skip branches where $|x_{pu}| <$ this threshold \\
\end{longtable}

Two preset builders are provided:
\begin{itemize}
  \item \texttt{make\_bp\_flags(min\_x=1e-6)} --- builds $\mathbf{B}'$ for the FDXB variant.
  \item \texttt{make\_bpp\_flags()} --- builds $\mathbf{B}''$ for the FDBX variant.
\end{itemize}

% ------------------------------------------------------------
\subsection{Projection Layer}
\label{sec:projection}
% ------------------------------------------------------------

The projection layer (\texttt{include/hacdcpf/projection/}) converts a
\emph{rich} \texttt{HybridPowerSystem} (with transformers, switches, complex
converters) into a \emph{canonical} form ready for solver assembly.

\subsubsection{Canonical Network Types}

Declared in \texttt{projection/canonical\_network.hpp}:

\begin{longtable}{ll}
\hline\hline
\textbf{Type} & \textbf{Description} \\
\hline\endhead\hline\endfoot
\texttt{BranchOriginType} & \texttt{Transformer2W}, \texttt{Transformer3W}, \texttt{Switch} \\
\texttt{BranchExpandEntry} & Records which branch was synthesised from which rich element \\
\texttt{BranchExpandMap} & Collection of \texttt{BranchExpandEntry} records \\
\hline
\texttt{BusMergeMap::ext\_to\_int} & External bus index $\to$ merged internal position \\
\texttt{BusMergeMap::int\_to\_ext} & Internal position $\to$ representative external index \\
\texttt{BusMergeMap::groups} & Per-representative list of merged external buses \\
\texttt{BusMergeMap::dead\_bus\_indices} & External indices of dead-island buses stripped \\
\texttt{BusMergeMap::branch\_orig\_to\_proj} & Original branch pos $\to$ projected pos (-1 = dead) \\
\texttt{BusMergeMap::has\_merges()} & True if any zero-impedance merging occurred \\
\texttt{BusMergeMap::has\_dead\_buses()} & True if any dead islands were stripped \\
\texttt{BusMergeMap::identity(n)} & Constructs no-op (identity) mapping for $n$ buses \\
\hline\hline
\end{longtable}

\subsubsection{Projection Functions}

Declared in \texttt{projection/project\_to\_canonical.hpp}:

\paragraph{Count queries (read-only)}
\begin{center}
\begin{tabular}{ll}
\hline\hline
\textbf{Function} & \textbf{Returns} \\
\hline
\texttt{n\_ac\_buses(sys)}, \texttt{n\_dc\_buses(sys)}, \texttt{n\_buses(sys)} & Bus counts \\
\texttt{n\_ac\_branches(sys)}, \texttt{n\_dc\_branches(sys)}, \texttt{n\_branches(sys)} & Branch counts \\
\texttt{n\_generators(sys)}, \texttt{n\_loads(sys)}, \texttt{n\_shunts(sys)} & Component counts \\
\texttt{n\_storage(sys)}, \texttt{n\_renewable\_gens(sys)}, \texttt{n\_converters(sys)} & Component counts \\
\hline\hline
\end{tabular}
\end{center}

\paragraph{Aggregation queries}
\begin{center}
\begin{tabular}{ll}
\hline\hline
\textbf{Function} & \textbf{Description} \\
\hline
\texttt{total\_gen\_capacity\_mw(sys)} & Sum of all generator $P_{max}$ \\
\texttt{total\_gen\_pmin\_mw(sys)} & Sum of all generator $P_{min}$ \\
\texttt{total\_load\_p\_mw(sys)} / \texttt{q\_mvar} & Total AC active / reactive load \\
\texttt{total\_renewable\_capacity\_mw(sys)} & Total renewable nameplate capacity \\
\texttt{total\_storage\_capacity\_mwh(sys)} & Total storage energy capacity \\
\texttt{total\_storage\_power\_mw(sys)} & Total storage peak power \\
\texttt{total\_system\_inertia\_mws(sys)} & Total rotational inertia \\
\texttt{total\_emission\_tco2\_h(sys)} & Total CO\textsubscript{2} emission rate \\
\hline\hline
\end{tabular}
\end{center}

\paragraph{Topology helpers}
\begin{itemize}
  \item \texttt{ac\_adjacency\_list(sys)} --- adjacency list of AC buses.
  \item \texttt{generators\_per\_bus(sys)} --- generator count per bus.
  \item \texttt{print\_summary(sys, os)} / \texttt{summary\_string(sys)} --- formatted system summary.
\end{itemize}

\paragraph{Projection pipeline}

\texttt{project\_to\_canonical\_models(sys)} performs:
\begin{enumerate}
  \item Expand \texttt{Transformer2W} / \texttt{Transformer3W} into equivalent \(\pi\)-model \texttt{ACBranch} entries.
  \item Expand closed \texttt{Switch} devices into zero-impedance branches.
  \item Populate \texttt{BranchExpandMap} for result un-projection.
\end{enumerate}

Two overloads accept either \texttt{const} reference (copy) or rvalue (move-efficient).

\paragraph{Zero-impedance bus merging}

\texttt{merge\_zero\_impedance\_buses(sys)} merges buses connected by branches
with $|Z| < \texttt{kBusMergeZThreshold}=10^{-4}$~pu into equivalence classes,
choosing one representative per class.  The resulting \texttt{BusMergeMap}
allows \texttt{unproject\_bus\_vector(merged, map)} to expand solver voltage
solutions back to the original bus indexing, returning 0.0 for dead-island
buses.

\texttt{strip\_dead\_islands(sys)} removes buses with no generation path
(neither direct generators nor grid connections).

% ------------------------------------------------------------
\subsection{Backend Solver Abstractions}
\label{sec:backend}
% ------------------------------------------------------------

The \texttt{include/hacdcpf/backend/} headers expose thin adapter namespaces
and abstract interfaces for the three solver categories used by the numerical
kernels.

\subsubsection{Linear Solver Backend}

\texttt{backend/linear\_solver.hpp} re-exports into \texttt{hacdcpf::backend}
and \texttt{hacdcpf::powerflow} (for backward compatibility):

\begin{itemize}
  \item \texttt{SparseLinearSolver} --- abstract base class for sparse LU-type solvers.
  \item \texttt{EigenSparseLUSolver} --- Eigen SparseLU concrete implementation.
  \item \texttt{make\_default\_sparse\_solver()} --- factory returning an \texttt{EigenSparseLUSolver}.
\end{itemize}

These are defined in \texttt{engine/kernel/linear\_algebra/linear\_solver.hpp}
and exposed here for layer-agnostic consumption.

\subsubsection{NLP Solver Backend}

\texttt{backend/nlp\_solver.hpp} --- namespace \texttt{hacdcpf::backend}.

\begin{longtable}{lll}
\hline\hline
\textbf{Type / Field} & \textbf{Default} & \textbf{Description} \\
\hline\endhead\hline\endfoot
\multicolumn{3}{l}{\textbf{NLPStatus}} \\
\texttt{Optimal} & --- & Global/local optimum found \\
\texttt{LocalOptimal} & --- & Local optimum, not proven global \\
\texttt{Infeasible} & --- & Problem is infeasible \\
\texttt{Unbounded} & --- & Objective is unbounded below \\
\texttt{MaxIterations} & --- & Iteration limit reached \\
\texttt{NumericalFailure} & --- & NaN/Inf or pivot breakdown \\
\hline
\multicolumn{3}{l}{\textbf{NLPResult}} \\
\texttt{status} & \texttt{Unknown} & Solve outcome \\
\texttt{x} & --- & Primal solution vector \\
\texttt{lambda} & --- & Equality constraint multipliers \\
\texttt{mu} & --- & Inequality constraint multipliers \\
\texttt{objective} & 0.0 & Optimal objective value \\
\texttt{iterations} & 0 & Inner iteration count \\
\texttt{primal\_infeasibility} & 0.0 & $\|\mathbf{g}(x^*)\|_\infty$ \\
\texttt{dual\_infeasibility} & 0.0 & $\|\nabla_x L\|_\infty$ \\
\texttt{complementarity} & 0.0 & KKT complementarity residual \\
\hline
\multicolumn{3}{l}{\textbf{NLPOptions}} \\
\texttt{feasibility\_tol} & \(10^{-6}\) & Constraint feasibility tolerance \\
\texttt{optimality\_tol} & \(10^{-6}\) & KKT optimality tolerance \\
\texttt{max\_iterations} & 300 & Maximum inner iterations \\
\texttt{verbose} & false & Print solver progress \\
\hline\hline
\end{tabular}

\subsubsection{QP Solver Backend}

\texttt{backend/qp\_solver.hpp} --- namespace \texttt{hacdcpf::backend}.

The \texttt{IQPSolver} abstract class solves:
\begin{equation}
  \min_x \;\tfrac{1}{2}x^\top H x + c^\top x \quad\text{s.t.}\quad Ax \leq b,\; l \leq x \leq u.
\end{equation}

\begin{longtable}{lll}
\hline\hline
\textbf{Type / Field} & \textbf{Default} & \textbf{Description} \\
\hline\endhead\hline\endfoot
\multicolumn{3}{l}{\textbf{QPStatus}} \\
\texttt{Optimal} & --- & Optimal solution found \\
\texttt{Infeasible, Unbounded} & --- & Problem degenerate \\
\texttt{MaxIterations, NumericalFailure} & --- & Solver limits \\
\hline
\multicolumn{3}{l}{\textbf{QPResult}} \\
\texttt{primal} & --- & Solution $x^*$ \\
\texttt{dual} & --- & Constraint multipliers \\
\texttt{objective} & 0.0 & Optimal objective \\
\texttt{iterations} & 0 & Simplex / interior-point iterations \\
\hline
\multicolumn{3}{l}{\textbf{QPOptions}} \\
\texttt{feasibility\_tol} & \(10^{-8}\) & Primal feasibility tolerance \\
\texttt{optimality\_tol} & \(10^{-8}\) & Dual optimality tolerance \\
\texttt{max\_iterations} & 10000 & Iteration cap \\
\hline\hline
\end{longtable}

% ------------------------------------------------------------
\subsection{Thread Pool}
\label{sec:thread-pool}
% ------------------------------------------------------------

\texttt{include/hacdcpf/util/thread\_pool.hpp} provides a lightweight fixed-size
worker pool for parallelising the AC Jacobian evaluation kernel.  It exposes
three core operations:

\begin{center}
\begin{tabular}{ll}
\hline\hline
\textbf{Member} & \textbf{Description} \\
\hline
\texttt{ThreadPool(num\_threads)} & Construct; 0 uses \texttt{hardware\_concurrency()} \\
\texttt{submit(callable)} & Enqueue a task; returns \texttt{std::future} \\
\texttt{parallel\_for(n, fn)} & Execute \texttt{fn(i)} for $i \in [0, n)$ across workers \\
\texttt{ThreadPool::global()} & Process-wide singleton (created on first call) \\
\hline\hline
\end{tabular}
\end{center}

The thread count is controlled by \texttt{PowerFlowOptions::ac\_eval\_threads};
setting it to~1 (the default) disables parallelism and uses the calling thread.
The pool is persistent across solve calls (no spin-up overhead per solve).

% ------------------------------------------------------------
\subsection{Detail Utilities}
\label{sec:detail}
% ------------------------------------------------------------

\texttt{include/hacdcpf/detail/} contains internal implementation helpers not
part of the public API:

\begin{center}
\begin{tabular}{ll}
\hline\hline
\textbf{Header} & \textbf{Contents} \\
\hline
\texttt{logging.hpp} & \texttt{hacdcpf::detail::log(level, msg)} debug logger \\
\texttt{string\_utils.hpp} & String trimming, splitting, and formatting helpers \\
\texttt{core\_compat.hpp} & C++17/20 portability shims (\texttt{std::span} polyfill, etc.) \\
\texttt{internal\_helpers.hpp} & Miscellaneous internal utilities (range checks, etc.) \\
\hline\hline
\end{tabular}
\end{center}

These headers are consumed only by \texttt{src/} translation units.  External
users should not include them directly; the public API encapsulates all
necessary behaviour.
